%% ch09 --- Trzy równania na pogodę: atraktor Lorenza
%%
%% Napisany we wrześniu 2026. Każda liczba, której czytelnik nie policzy w
%% pamięci, pochodzi z code/measure/ch09.py, który importuje TE SAME funkcje,
%% które drukują listingi (code/ch09/). Reguły Lorenza używają tylko + - * /,
%% więc każda trajektoria na tych stronach jest taka sama na każdej maszynie;
%% wykładnik jest zapisany z dwoma miejscami po przecinku, których ostatni
%% bit logarytmu nie ruszy.
%%
%% CZEGO TEN ROZDZIAŁ NIE MOŻE POWIEDZIEĆ. Słowo "dziwny" i geometria
%% atraktora należą do rozdziału 10, mierzenie takich kształtów do 11,
%% pogoda kontra klimat i prognozy zespołowe do 14, zaufanie do statystyk
%% policzonych komputerem do 16, synchronizacja do 17, a historia z 1961 i
%% motyl z 1972 do rozdziału 6.
%%
%% Bliźniak: chapters/en/ch09-lorenz.tex. Podrozdział za podrozdziałem, makro
%% za makrem, liczba za liczbą; tools/parity.py je porównuje.

\chapter{Trzy równania na pogodę: atraktor Lorenza}
\label{ch:lorenz}
\index{układ Lorenza}\index{Lorenz, Edward}

\noindent
Postaw szeroki garnek z wodą na małym ogniu. Ciepła woda przy dnie chce się
unieść, a chłodna przy powierzchni opaść, i powyżej pewnej temperatury woda
zaczyna się toczyć: w górę z jednej strony, w dół z drugiej, jak powolne
koło. Powietrze robi to samo nad nasłonecznionym polem, a w większej skali to
toczenie się jest jednym z silników pogody.

Opisz ten wir trzema liczbami i trzema krótkimi regułami, dokładnymi, bez
żadnej przypadkowości. Czy będzie się kręcił równo, jak koło? Czy też trzy
liczby mogą być tak nieregularne jak niebo? Część~II znalazła chaos w
regułach, które robią jeden krok na raz. Ten rozdział znajduje go w czymś,
co płynie \dash{} a potem znajduje wewnątrz chaosu coś, czego nikt się nie
spodziewał: kształt.

\begin{outcomes}
\outcome{przeczytać trzy reguły zmiany Lorenza wyraz po wyrazie i powiedzieć,
  co każdy wyraz robi z wirem}
\outcome{prowadzić reguły metodą RK4 i odczytać wynik z dwóch widoków: jedna
  liczba względem drugiej oraz jedna liczba względem czasu}
\outcome{zmierzyć, jak szybko rozchodzą się dwa prawie identyczne przebiegi,
  i zamienić to tempo na horyzont przewidywania}
\outcome{powiedzieć, czym jest atraktor i dlaczego układ chaotyczny nie może
  pójść, gdzie zechce}
\outcome{policzyć statystykę przebiegu i sprawdzić, że nie zależy od tego,
  skąd przebieg wystartował}
\end{outcomes}

\section{Trzy liczby na garnek wody}
\label{sec:ch09-rules}
\index{konwekcja}

\noindent
W 1963 roku Edward Lorenz, meteorolog z Massachusetts Institute of
Technology, opublikował pracę zatytułowaną \enquote{Deterministic Nonperiodic
Flow}, czyli \enquote{Deterministyczny nieokresowy przepływ}. Ten tytuł to
cały rozdział w pigułce. \emph{Deterministyczny}: w regułach nie ma
przypadku. \emph{Nieokresowy}: ruch nigdy się nie powtarza.
\emph{Przepływ}: czas płynie gładko, więc reguła podaje prędkość, z jaką
zmienia się stan, a nie jego następną wartość, jak w
rozdziale~\ref{ch:motion}. Model Lorenza to drastycznie okrojony opis płynu
podgrzewanego od dołu i zachowuje trzy liczby:

\begin{itemize}
\item $x$: jak szybko kręci się wir i w którą stronę (dodatnie w jedną,
  ujemne w drugą);
\item $y$: o ile cieplejsza jest strona, którą woda się wznosi, od strony,
  którą opada;
\item $z$: o ile temperatura, od dołu do góry, została odgięta od
  równomiernego spadku przez mieszanie, które robi wir.
\end{itemize}

\noindent
Zapisuj $\mathrm{d}x/\mathrm{d}t$ i czytaj \enquote{tempo $x$}: to, jak
szybko $x$ zmienia się w tej chwili; to nazwa prędkości, a nie ułamek do
skracania. Reguły Lorenza to
\[
\frac{\mathrm{d}x}{\mathrm{d}t} = \sigma\,(y - x), \qquad
\frac{\mathrm{d}y}{\mathrm{d}t} = x\,(\rho - z) - y, \qquad
\frac{\mathrm{d}z}{\mathrm{d}t} = x\,y - \beta\,z,
\]
z trzema stałymi liczbami nazwanymi greckimi literami: $\sigma$ (sigma) to
$10$, $\rho$ (ro) to $28$, a $\beta$ (beta) to $\tfrac{8}{3}$. Słowami:

\begin{itemize}
\item \textbf{Wir goni różnicę temperatur.} Gdy $y$ jest większe od $x$, wir
  przyspiesza, a gdy mniejsze, zwalnia; $\sigma$ mówi, jak żwawo.
\item \textbf{Wir buduje różnicę, a różnica ucieka.} Po wymnożeniu druga
  reguła to $\rho x - xz - y$. Kręcący się wir wynosi ciepłą wodę w górę z
  jednej strony i sprowadza chłodną w dół z drugiej, budując $y$ w tempie
  $\rho x$, gdzie $\rho$ mówi, jak mocno grzejemy garnek ($\rho = 1$ to
  najsłabsze grzanie, przy którym wir w ogóle rusza). Wyraz $-y$ to różnica
  uciekająca na boki; wyraz $-xz$ to mieszanie, które się odgryza.
\item \textbf{Wir miesza warstwę, a warstwa się uspokaja.} Wir wynoszący
  ciepłą wodę w górę ($x$ i $y$ tego samego znaku) wygina profil temperatury w
  tempie $xy$; pozostawiony sam sobie profil wraca w tempie $\beta z$, gdzie
  $\beta$ wynika z kształtu wirów.
\end{itemize}

\noindent
Tylko dwa wyrazy, $xz$ i $xy$, mnożą przez siebie dwie z trzech liczb; każdy
inny wyraz to stała liczba razy jedna z nich. Te dwa iloczyny to cała
nieliniowość modelu w sensie z rozdziału~\ref{ch:feedback} i okaże się, że
wystarczą.

\mermaidfig{ch09-feedback}{Trzy liczby popychają się nawzajem w kółko. Wir
buduje różnicę temperatur, różnica napędza wir, a razem mieszają warstwę,
która osłabia różnicę.}{trzy liczby w pętli}

\begin{think}
Załóżmy, że woda stoi zupełnie nieruchomo: $x = 0$, $y = 0$ i $z = 0$. Co
mówią trzy reguły?
\end{think}
\reveal{Każde tempo jest zerem: $10 \times (0 - 0) = 0$, potem $0 - 0 = 0$,
potem $0 - 0 = 0$. Nic się nie zmienia.}
To garnek, który po cichu przekazuje ciepło w górę, bez żadnego wiru. Stan, w
którym każde tempo jest zerem, to \emph{punkt stały} przepływu, w języku
rozdziału~\ref{ch:iteration}: zacznij w nim, a w nim zostaniesz. Nie jest
jedyny.

\begin{think}
Sprawdź $x = y = \sqrt{72}$, czyli około \val{ch09.fixed.xy}, oraz $z = 27$.
Czy wszystkie trzy tempa są zerem? Potrzebujesz tylko
$\tfrac{8}{3} \times 27 = 72$.
\end{think}
\reveal{Tak. Pierwsze tempo to $10 \times 0 = 0$. Drugie to
$x \times (28 - 27) - y = x - y = 0$. Trzecie to $x \times x - 72 = 72 -
72 = 0$.}
To \emph{ustalony wir}: kręci się ze stałą prędkością przy stałej różnicy
temperatur. Z minusami, $x = y = -\sqrt{72}$ i $z = 27$, to ten sam wir
kręcący się w drugą stronę. Dla każdego grzania $\rho$ większego od $1$ oba
ustalone wiry leżą w $x = y = \pm\sqrt{\beta(\rho - 1)}$ i $z = \rho - 1$,
gdzie $\pm$ znaczy \enquote{plus albo minus}.

\begin{lab}{l09_01_rates}{Reguły Lorenza własną ręką}
Napisz \code{rates}, która dostaje stan $(x, y, z)$ i zwraca trzy tempa,
oraz \code{fixed\_points}, która zwraca każdy stan, w którym wszystkie trzy
są zerem. Test porównuje twoje tempa z biblioteką książki i sprawdza, że
naprawdę są zerem w każdym punkcie, który zwrócisz.
\end{lab}

\section{Za przepływem}
\label{sec:ch09-flow}
\index{układ Lorenza!motyl}\index{RK4}

\noindent
Oto te same reguły w wersji z biblioteki książki. Ostatnie trzy wiersze to
trzy reguły; \api{lorenz} zwraca funkcję, która dostaje stan i oddaje jego
trzy tempa, czyli dokładnie to, o co prosiło twoje laboratorium.

\pyregion{code/src/chaoslab/systems.py}{lorenz}{Reguły Lorenza z 1963 roku w chaoslab}{lst:ch09-field}

\noindent
Żeby za nimi podążać, użyj metody RK4 z rozdziału~\ref{ch:motion}:
\api{integrate} robi kroki długości \code{DT}, czyli jednej setnej jednostki
czasu, i zapamiętuje każdy stan. Jednostka czasu jest własną jednostką
modelu, a nie sekundą ani godziną.

\pyregion{code/ch09/lorenz_run.py}{run}{Czterdzieści jednostek czasu przepływu Lorenza, od stanu (1, 1, 1)}{lst:ch09-run}

\begin{think}
Wahadło z tarciem wchodzi spiralą w swój punkt spoczynku. Garnek Lorenza ma
trzy punkty stałe. Czy przebieg rozpoczęty w $(1, 1, 1)$, blisko
nieruchomej wody, osiądzie w jednym z nich?
\end{think}
\reveal{Nie. Zostaje odepchnięty od nieruchomej wody, krąży wokół jednego
ustalonego wiru coraz szerszymi pętlami, zostaje przerzucony, by krążyć
wokół drugiego, i nigdy się nie uspokaja.}
Przy tym grzaniu żaden ustalony wir nie potrafi zatrzymać stanu; potrafią
tylko wtedy, gdy $\rho$ jest mniejsze od $\tfrac{470}{19}$, czyli około
\val{ch09.rho.hopf} \dash{} tę liczbę daje z reguł algebra, której ta
książka nie przeprowadza. Zmniejsz grzanie do $20$, a ten sam start
rzeczywiście się uspokoi: ostatni wiersz wydruku to ustalony wir w
$x = y = -\val{ch09.steady.x}$, $z = 19$.

\transcript{ch09-butterfly}

\noindent
Wiersz oznaczony \code{loops} wypisuje jedną literę na pętlę: \code{L} dla
pętli po stronie, gdzie $x$ jest ujemne, \code{R} dla drugiej. Przebieg
robi \val{ch09.loops.left} pętli po lewej, zanim zrobi pierwszą po prawej,
a potem litery idą seriami od jednej do pięciu, bez żadnego wzoru. Nie ma go
czego szukać: to jest \emph{nieokresowość} z tytułu Lorenza.

\plotfig{ch09-butterfly}{Przebieg od $(1, 1, 1)$. Z lewej: $x$ względem $z$
od chwili 5 do chwili 60 \dash{} pętle wokół dwóch ustalonych wirów
(krzyżyki) rysują dwa skrzydła. Z prawej: $x$ względem czasu przez
pierwsze czterdzieści jednostek czasu; linie przerywane to ustalone
wiry.}{motyl z dwóch stron}

\noindent
Lewa połowa rysunku~\ref{fig:ch09-butterfly} zasłużyła na nazwę
\emph{motyla Lorenza}. Każde skrzydło to wir kręcący się w jedną stronę;
przejście między nimi to wir zmieniający kierunek.

\begin{think}
Służby pogodowe liczą dziś modele niosące wiele milionów liczb. Model
Lorenza niesie trzy. Czy trzy wystarczą na prawdziwy chaos, czy tylko na
jego szkic?
\end{think}
\reveal{Trzy wystarczą, i trzy to najmniej, ile potrzebuje przepływ.}
Rozdział~\ref{ch:motion} podał tę dolną granicę; oto powód. Na płaszczyźnie
droga nigdy nie może przeciąć samej siebie, bo dwie drogi przez jeden punkt
miałyby tę samą przyszłość; przepływ z dwiema liczbami może więc tylko
osiąść, zwijać się spiralą albo krążyć po stałej pętli. Z trzecią liczbą
droga może przechodzić nad sobą i pod sobą, i to wystarcza.

\begin{trapbox}
\textbf{\enquote{Chaos potrzebuje wielu zmiennych.}} Potrzebuje trzech liczb
i dwóch mnożeń. Pogodę trudno przewidzieć nie tylko dlatego, że atmosfera
jest wielka: garnek wody okrojony do trzech liczb jest już
nieprzewidywalny. Czego chaos naprawdę potrzebuje, to temat
rozdziału~\ref{ch:strange}.
\end{trapbox}

\begin{worldbox}
\emph{Chaotyczne koło wodne} to garnek Lorenza zbudowany z drewna i wody:
cieknące kubki wiszą na obwodzie koła, a woda wlewa się do kubków na górze.
Willem Malkus i Louis Howard zbudowali takie koło w MIT, a jego równania, po
przemianowaniu wielkości, są równaniami Lorenza: koło kręci się, zwalnia i
zawraca w chwilach, których nikt nie potrafi zapowiedzieć. Równania pewnego
rodzaju lasera też są równaniami Lorenza, co pokazał Hermann Haken w 1975
roku.
\end{worldbox}

\section{Dwa starty o włos od siebie}
\label{sec:ch09-apart}
\index{wrażliwość na warunki początkowe}\index{wykładnik Lapunowa!przepływu}

\noindent
Rozdział~\ref{ch:butterfly} zmierzył wrażliwość na odwzorowaniu. Oto ten sam
pomiar na przepływie. Przebieg najpierw idzie przez pięćdziesiąt jednostek
czasu, żeby zostawić start za sobą; potem zostaje skopiowany, a kopia
przesunięta o jedną stumilionową w $x$.

\pyregion{code/ch09/two_runs.py}{split}{Dwa przebiegi o jedną stumilionową od siebie}{lst:ch09-split}

\noindent
Oba są prowadzone tą samą regułą, a program co dwie jednostki czasu wypisuje
$x$ każdego z nich i odległość między nimi:

\transcript{ch09-two-runs}

\noindent
Do chwili 14 oba przebiegi zgadzają się co do dwóch miejsc po przecinku; w
chwili 24 są na różnych skrzydłach. Odległość rośnie mniej więcej
dziesięciokrotnie co dwie, trzy jednostki czasu, choć nie gładko: puchnie i
kurczy się, gdy para obiega pętlę. Pierwszy raz przekracza $1$ w chwili
\val{ch09.apart.first}.

Rozdział~\ref{ch:lyapunov} nazwał to tempo wykładnikiem Lapunowa, $\lambda$
(lambda): to średnia logarytmu tego, jak bardzo rozciąga się mała szczelina.
Dla przepływu średnia liczona jest na jednostkę czasu. Biblioteka mierzy go
tak:

\pyregion{code/src/chaoslab/measure.py}{flowexp}{Wykładnik Lapunowa przepływu}{lst:ch09-flowexp}

\noindent
Prowadź dwa przebiegi o włos od siebie. Po każdym kroku dodaj logarytm tego,
o ile urosła szczelina między nimi, a potem ściągnij drugi przebieg z
powrotem na pierwotną małą odległość, w tym samym kierunku, żeby szczelina
nigdy nie urosła do rozmiarów motyla, gdzie przestałaby mierzyć rozciąganie.
Suma podzielona przez upływ czasu to $\lambda$. Oto on z trzech startów:

\pyregion{code/ch09/exponent.py}{exponent}{Wykładnik przepływu Lorenza z trzech startów}{lst:ch09-exponent}

\transcript{ch09-exponent}

\noindent
Od \val{ch09.lambda.lo} do \val{ch09.lambda.hi} na jednostkę czasu, z
trzech zupełnie różnych startów: przyjmijmy \val{ch09.lambda}. Średnio
szczelina mnoży się przez $e$ do potęgi \val{ch09.lambda}, czyli około
\val{ch09.stretch}, na każdą jednostkę czasu.

\begin{think}
Horyzont z rozdziału~\ref{ch:lyapunov}: błąd rośnie do tolerancji w czasie
około $\ln(\text{tolerancja} / \text{błąd}) / \lambda$. Przy błędzie
$10^{-8}$, tolerancji $1$ i $\lambda = \val{ch09.lambda}$: jak długo dwa
przebiegi powinny pozostać bliżej niż $1$ od siebie? Potrzebujesz
$\ln(10^{8})$, czyli około \val{ch09.ln.eight}.
\end{think}
\reveal{Około $\val{ch09.ln.eight} / \val{ch09.lambda}$, czyli mniej więcej
\val{ch09.horizon} jednostek czasu. Para powyżej potrzebowała
\val{ch09.apart.first}.}
Jedna para to jedna próbka. Rozdziel przebieg w \val{ch09.apart.pairs}
różnych miejscach wzdłuż motyla, a pierwsze przekroczenie $1$ wypada
gdziekolwiek od chwili \val{ch09.apart.min} do chwili \val{ch09.apart.max},
średnio w \val{ch09.apart.mean}. Wzór przewiduje średnią; rozrzut bierze się
stąd, że rozciąganie nie jest takie samo na całych skrzydłach.

\plotfig{ch09-apart}{Dwa przebiegi, które startują o jedną stumilionową od
siebie. Z lewej: ich $x$ zgadzają się, a potem się rozchodzą. Z prawej: ich
odległość w skali logarytmicznej wspina się wzdłuż przerywanej linii
wykładnika, aż osiągnie rozmiar motyla.}{dwa rozchodzące się przebiegi}

\noindent
W regule nie ma nic niepewnego, a mimo to prognoza od startu znanego z
dokładnością do ośmiu miejsc po przecinku wystarcza na mniej więcej
dwadzieścia jednostek czasu. (Dwa przebiegi, które się rozeszły, można
znowu zmusić do zgody, pozwalając jednemu czuć drugi: to opowieść
rozdziału~\ref{ch:taming}.)

\section{Nieprzewidywalny, ale nie gdziekolwiek}
\label{sec:ch09-attractor}
\index{atraktor}

\noindent
Jeśli przebiegu nie da się prognozować dalej niż na dwadzieścia jednostek
czasu, co w ogóle można o nim powiedzieć?

\begin{think}
W chwili 24 dwa przebiegi powyżej nie mają ze sobą nic wspólnego: jeden ma
$x$ w pobliżu $-12$, drugi w pobliżu $8$. Puść je jeszcze na tysiąc jednostek
czasu. Czy któryś z nich może się znaleźć w $x = 100$ albo w $z = -50$?
\end{think}
\reveal{Nie. Przez \val{ch09.long.time} jednostek czasu $x$ nie oddala się od
zera bardziej niż o \val{ch09.long.x} w żadną stronę, a $z$ pozostaje między
\val{ch09.long.zlo} a \val{ch09.long.zhi}. I tak jest niezależnie od tego,
skąd wystartujesz.}
Sprawdź to \enquote{niezależnie od tego, skąd} na ośmiu startach w
narożnikach pudła o boku osiemdziesiąt, daleko poza wszystkim, co widać na
rysunkach powyżej:

\pyregion{code/ch09/cloud.py}{cloud}{Osiem startów daleko poza motylem}{lst:ch09-cloud}

\noindent
Program prowadzi każdy przez pięćdziesiąt jednostek czasu, pomija pierwsze
pięć i wypisuje pokój, w którym żył każdy przebieg:

\transcript{ch09-cloud}

\noindent
Każdy przebieg zostaje ściągnięty ze swojego narożnika i potem mieszka w tym
samym pokoju co długi przebieg. Zauważ też, że wiersze idą w identycznych
parach: $(-40, -40, -40)$ i $(40, 40, -40)$ dają te same liczby. Następny
podrozdział mówi dlaczego.

\begin{trapbox}
\textbf{\enquote{Układ chaotyczny może pójść gdziekolwiek.}}
Nieprzewidywalny to nie to samo co nieograniczony. Nie powiesz, w którym
miejscu motyla będzie przebieg Lorenza za sto jednostek czasu, ale powiesz,
że będzie na motylu, i to jest prawdziwe przewidywanie: cienki zbiór stanów
spośród wszystkich, jakie mogłyby przyjąć trzy liczby.
\end{trapbox}

\noindent
Zbiór stanów, na który przebiegi zostają ściągnięte i którego potem nigdy
nie opuszczają, to \emph{atraktor}. Wahadło z tarciem ma za atraktor jeden
punkt; wahadło zegara, podtrzymywane w ruchu, ma pętlę. Atraktorem Lorenza
jest motyl, ten sam z niemal każdego startu, a przebieg na nim nigdy się
nie powtarza i nigdy go nie opuszcza. Jak droga może bez końca krążyć w
ograniczonym obszarze, nie przecinając się i nie wracając po własnych
śladach, to zagadka dotycząca samego kształtu \dash{} dlatego nazywa się go
\emph{atraktorem dziwnym}. Rozdział~\ref{ch:strange} otwiera ten kształt, a
rozdział~\ref{ch:fractals} mierzy kształty takie jak on.

\begin{rigourbox}
Dwie rzeczy są tu pokazane przez uruchomienie, a nie udowodnione. To, że
żaden przebieg nie ucieknie, da się udowodnić na stronie algebry o
wielkości, która maleje zawsze, gdy stan jest daleko. To, że motyl
narysowany przez komputer naprawdę istnieje i naprawdę jest chaotyczny, a
nie jest wytworem kroków i zaokrągleń, pozostawało otwarte przez ponad
trzydzieści lat, aż Warwick Tucker rozstrzygnął to dowodem, w którym
komputer liczył tak, że każdy błąd zaokrąglenia był wzięty pod uwagę.
\end{rigourbox}

\section{Co wciąż da się przewidzieć}
\label{sec:ch09-statistics}
\index{atraktor!statystyki}

\noindent
Motyl mówi, gdzie przebieg mieszka. Następne pytanie brzmi: jak spędza tam
czas.

\begin{think}
Nie powiesz, na którym skrzydle będzie przebieg w chwili 50. Czy powiesz,
jaką część długiego przebiegu spędza on na prawym skrzydle, gdzie $x$ jest
dodatnie?
\end{think}
\reveal{Tak: połowę. Zmień jednocześnie znaki $x$ i $y$: pierwsze dwa tempa
zmieniają znak razem z nimi, a trzecie, $xy - \beta z$, wcale się nie
zmienia, bo $(-x)(-y) = xy$.}
Reguły nie dbają więc o to, w którą stronę kręci się wir. Każdy przebieg na
jednym skrzydle ma lustrzane odbicie na drugim \dash{} dlatego wiersze
chmury szły parami \dash{} i długi przebieg dzieli swój czas między skrzydła
po równo.

\begin{trapbox}
\textbf{\enquote{Jeśli droga jest nieprzewidywalna, nic w niej nie da się
przewidzieć.}} Droga jest nieprzewidywalna dalej niż na mniej więcej
dwadzieścia jednostek czasu. Kształt, na którym mieszka, już nie, podobnie
jak to, ile czasu spędza w każdej jego części. To są \emph{statystyki
atraktora} i z każdego startu wychodzą takie same.
\end{trapbox}

\begin{lab}{l09_02_wings}{Dwie liczby, o których start nie decyduje}
Napisz \code{fraction\_right}, czyli udział stanów przebiegu z $x$ większym od
zera, oraz \code{mean\_height}, czyli średnią jego $z$. Test prowadzi reguły
Lorenza od $(1, 1, 1)$ i od dalekiego narożnika chmury i sprawdza, że twoje
dwa podsumowania zgadzają się między przebiegami.
\end{lab}

\noindent
Oto wersje z książki:

\pyregion{code/ch09/averages.py}{stats}{Dwie statystyki przebiegu}{lst:ch09-stats}

\noindent
Listing prowadzi pięć startów, każdy przez 500 jednostek czasu, odrzuca
pierwsze dziesięć, żeby start został zapomniany, i dla każdego przebiegu
wypisuje, na którym skrzydle jest w chwili 50, jaką część czasu spędził na
prawym skrzydle i jaka jest jego średnia wysokość:

\transcript{ch09-averages}

\noindent
Środkowa kolumna to część nieprzewidywalna: lewe, lewe, prawe, lewe, prawe, i
nic, co mógłbyś wiedzieć zawczasu, tego nie wybiera. Ostatnie dwie kolumny to
część przewidywalna. Udział prawego skrzydła wynosi od \val{ch09.frac.lo} do
\val{ch09.frac.hi} z każdego startu, tak jak zapowiadało lustro. Średnia
wysokość wynosi od \val{ch09.height.lo} do \val{ch09.height.hi}, a tej liczby
nie podpowiedziała żadna symetria: po prostu jest taka sama z każdego startu.

Statystyka skończonego przebiegu trochę się chwieje. Potnij jeden przebieg
długości \val{ch09.long.time} jednostek czasu na kawałki: w kawałkach po 100
jednostek czasu udział prawego skrzydła odbiega zwykle od połowy o około
\val{ch09.wobble.short}; w kawałkach po 2000 \dash{} o około
\val{ch09.wobble.long}; w całym przebiegu wynosi \val{ch09.long.frac}. Im
dłużej patrzysz, tym ostrzejsza statystyka \dash{} odwrotnie niż droga,
której prognoza jest tym gorsza, im dalej sięgasz.

\plotfig{ch09-histograms}{Ile czasu przebieg spędza przy każdej wartości
$x$ i $z$, z trzech startów, po 500 jednostek czasu każdy. Trzy obrysy leżą
jeden na drugim: histogram nie dba o to, skąd przebieg
wystartował.}{histogramy obojętne na start}

\noindent
Cały histogram tego, gdzie przebieg spędza czas, wychodzi taki sam z trzech
startów. To oś tego rozdziału i druga wielka myśl tej książki:
\textbf{indywidualna przyszłość układu chaotycznego jest nieprzewidywalna,
ale to, gdzie układ mieszka i jak często odwiedza każdą część, da się
przewidzieć}. Nawet $\lambda$ jest liczbą tego rodzaju.

Rozdział~\ref{ch:weather} buduje właśnie na tym rozróżnienie między pytaniem
o pogodę na przyszły tydzień a pytaniem o klimat i pokazuje, jak synoptycy
liczą wiele kopii naraz. Rozdział~\ref{ch:computers} pyta, dlaczego można
ufać statystyce, choć własne zaokrąglenia komputera szybko czynią jego drogę
inną niż prawdziwa.

\begin{summarybox}
\begin{itemize}
\item Model Lorenza z 1963 roku, opisujący podgrzewany garnek, zachowuje trzy
  liczby: prędkość wiru $x$, różnicę temperatur $y$ i wymieszanie $z$. Tylko
  dwa wyrazy jego trzech \textbf{reguł zmiany}, $xy$ i $xz$, są
  \textbf{nieliniowe}. Jego \textbf{punkty stałe} to nieruchomy garnek i dwa
  \textbf{ustalone wiry}.
\item Przy $\rho = 28$ przebieg nie osiada w żadnym punkcie stałym: krąży
  wokół jednego ustalonego wiru, potem wokół drugiego, przeskakując
  nieregularnie i nigdy się nie powtarzając. Trzy liczby to najmniej, ile
  \textbf{przepływ} potrzebuje do chaosu, i wystarczą.
\item Dwa przebiegi o jedną stumilionową od siebie rozchodzą się w mniej
  więcej dwadzieścia jednostek czasu. \textbf{Wykładnik Lapunowa} przepływu
  to około \val{ch09.lambda} na jednostkę czasu, a wzór na horyzont z
  rozdziału~\ref{ch:lyapunov} przewiduje średni czas rozejścia.
\item Każdy przebieg zostaje ściągnięty na ten sam ograniczony zbiór w
  kształcie motyla: \textbf{atraktor}. Układ chaotyczny nie może pójść, gdzie
  zechce.
\item \textbf{Statystyki atraktora}, takie jak część czasu na każdym
  skrzydle i średnia wysokość, są takie same z każdego startu. Droga jest
  nieprzewidywalna; to, gdzie układ mieszka, nie.
\end{itemize}
\end{summarybox}

\canyou

\begin{checkpoint}
\item Co słowami robi z wirem reguła na tempo $x$, $\sigma(y - x)$, gdy $y$
  jest większe od $x$?
  \answerto{Sprawia, że $x$ rośnie: wir przyspiesza w stronę różnicy
  temperatur.}
\item Gdzie przy $\rho = 28$ leżą trzy punkty stałe reguł Lorenza?
  \answerto{Nieruchomy garnek, $(0, 0, 0)$, oraz dwa ustalone wiry,
  $x = y = \pm\sqrt{72}$ (około \val{ch09.fixed.xy}) z $z = 27$.}
\item Dlaczego przepływ z dwiema liczbami nie może być chaotyczny i ile liczb
  używa model Lorenza?
  \answerto{Na płaszczyźnie droga nie może przeciąć samej siebie, więc może
  tylko osiąść, zwijać się spiralą albo krążyć po stałej pętli. Model
  Lorenza używa trzech, najmniej, ile pozwala drodze przechodzić nad sobą i
  pod sobą.}
\item Wykładnik to około \val{ch09.lambda} na jednostkę czasu. Gdyby start
  był znany tysiąc razy dokładniej, o ile dłużej dwa przebiegi pozostałyby
  bliżej niż $1$ od siebie?
  \answerto{O około $\ln(1000) / \val{ch09.lambda}$, czyli mniej więcej
  \val{ch09.gain.thousand} jednostek czasu: z mniej więcej dwudziestu do
  mniej więcej dwudziestu ośmiu, a nie tysiąc razy dłużej.}
\item Ktoś twierdzi, że przebieg Lorenza, skoro jest chaotyczny, może
  kiedyś dojść do $z = 200$. Co mu odpowiesz?
  \answerto{Nie dojdzie: każdy przebieg zostaje ściągnięty na ten sam
  ograniczony atraktor, na którym $z$ pozostaje mniej więcej między
  \val{ch09.long.zlo} a \val{ch09.long.zhi}.}
\item Jaką część długiego przebiegu zajmuje prawe skrzydło i skąd możesz to
  wiedzieć, niczego nie uruchamiając?
  \answerto{Połowę. Jednoczesna zmiana znaków $x$ i $y$ nie zmienia reguł,
  więc oba skrzydła są swoimi lustrzanymi odbiciami.}
\item Podaj liczbę opisującą przebieg Lorenza, której nie przewiduje żadna
  symetria, a która z każdego startu wychodzi taka sama.
  \answerto{Średnia wysokość $z$, od \val{ch09.height.lo} do
  \val{ch09.height.hi} z pięciu różnych startów; wykładnik Lapunowa, około
  \val{ch09.lambda}, to kolejna.}
\end{checkpoint}
